\section{Introduction}
The number of exactly distinguishable posterior directions can jump when
an additive relation between observation exponents disappears. The
persistent memory needed at a fixed prediction accuracy need not jump:
new directions may have arbitrarily small amplitude. We study this
question for positive sparse monomial experiments, and ask in addition
whether their sharp memory order can be realized by a machine whose
program contains only finitely many numerical bits.

Let $A=\{0=a_0<a_1<\cdots<a_{r-1}=D\}$, let
$W_A=\operatorname{span}\{t^a:a\in A\}$, and let $mA$ be the
set of sums of exactly $m$ members of $A$, with $0A=\{0\}$.
Write $h_A(m)=|mA|$. For the exact information statements the parameter
interval is $I=[l,u]$, $0\le l<u<\infty$, and the prior $\mu$ has
full support. For collision-uniform statements we take $I=[0,1]$.
A positive finite detector spans $W_A$, and a command in
$[\eta,1-\eta]^J$ retains or rejects each detector symbol. Every
reported symbol, including rejection, is counted as a trial. Its
likelihood is bounded below by a fixed positive number $\kappa$.
The precise experiment and common future-only equivalence are defined
in Section~\ref{sec:experiment}.

The one-step exponents range over a fixed compact subset $K$ of the
strictly ordered positive chamber. One fixed full-column-rank coefficient
matrix specifies the detector on $K$, with uniform positivity. Neither
semialgebraicity of $K$ nor a density for $\mu$ is assumed. The horizon
$N\ge2$ and prior are fixed. Choose a fixed
$H\ge\max\{1,N\max_Ka_{r-1}\}$. For each future length $m$, retain
all nonconstant formal homogeneous labels $\alpha\in\mathbb N_0^r$,
$|\alpha|=m$, including labels that have the same exponent. Their
normalized nodes are $x_\alpha=H^{-1}\alpha\cdot a$ and their number
is $q_m=\binom{m+r-1}{r-1}-1$. Set
\[
 \mathcal V_{m,\ell}(a)=\max_{|F|=\ell}
         \prod_{\{\alpha,\beta\}\subset F}|x_\alpha-x_\beta|,
 \qquad \mathcal V_{m,1}=1,
\]
where $F$ ranges over subsets of formal labels. Zero volumes remain
zero. Put $p_{n,m}=\min\{n(r-1),q_m\}$ and
\[
 \Psi_{n,m}(M,a)=\max_{1\le\ell\le p_{n,m}}
           (\mathcal V_{m,\ell}(a)/M)^{2/\ell},\qquad
 \Xi_N(M,a)=\max_{1\le n<N}\Psi_{n,N-n}(M,a).
\]
The checkpoint profile is zero at an endpoint. At each checkpoint the
query is selected independently after compression from a fixed finite
spanning menu of $m$-trial product probes. The excess squared loss is
the average squared error of its conditional event probabilities.
The notation $\overline R_{M;n,m}^a$ denotes optimal unconditional
checkpoint regret under independent uniform acquisition commands;
$R_{M;n,m}^{a,\max}$ denotes optimal worst-history regret.
Their causal counterparts $\mathfrak R^{a,\mathrm{av}}_{M,N}$ and
$\mathfrak R^{a,\max}_{M,N}$ minimize the maximum checkpoint criterion
over one stage-compatible filter. Only one checkpoint and its remaining
query trials are executed in a run.

\begin{theorem}[Finite numerical realization of sharp causal memory]
\label{thm:effective-main}
Under the preceding fixed-experiment hypotheses, the exact-real
checkpoint and causal laws have orders $\Psi_{n,m}$ and $\Xi_N$,
respectively, uniformly in $a\in K$ and integers $M\ge1$.
If the command cube is rationally specified and finite approximations
to the calibration coefficients and prior moments through formal degree
$N$, together with certified coefficient and positivity bounds, are supplied, the causal upper law is realized by integer transition
tables and dyadic output tables. With command and offline evaluation
errors at most $c/(M+1)$, one such table system has at most $M$
persistent states per stage and regret at most $C\Xi_N(M,a)$ on
every history. With a fixed memoryless command-coding rule, every such
system has maximum unconditional checkpoint
regret at least $c'\Xi_N(M,a)$ under the same exploration law.

The required numerical precision is $\log_2(M+1)+O(1)$ places and the
read-only program has $O(M^{J+1}\log(M+1))$ bits for the fixed
experiment. With computable numerical input, the tables are produced
by one uniform terminating algorithm. No online real-valued table,
posterior vector or discarded prefix is stored. No exponent-gap inverse
or decision of exact additive collision is used. The constants may
depend on the fixed prior and horizon; they do not depend on the
calibration or the state budget.
\end{theorem}

The exact-real part is established by the attainable pairing and intrinsic
classification below. The finite-table assertion is
Theorem~\ref{thm:digital-intrinsic}. For arbitrary, possibly noncomputable
full-support priors the mathematical classification is unchanged; an
algorithm for obtaining unsupplied noncomputable moments is not asserted.
The program size, numerical input precision and persistent state are
three different resources. With an external clock the persistent state
has $M$ values; storing the clock as well costs a factor depending only
on $N$. The program may be much larger than its changing state.

\subsection{Structure of the argument}
The geometric step is a pairing between an attainable product tangent
and a complete future test flag. At one interior binomial tuple the
product tangent contains $n(r-1)+1$ distinct monomials. Strict mixed
moment positivity holds for every full-support prior. Complete Hermite
blocks retain this positivity when nodes coincide; normalization removes
exactly one rank. The resulting local image carries an unconditional
volume minorization with the probability of acquiring the failure prefix
included. A separate bounded-format estimate covers the entire reachable
image, with dimension truncated at $p_{n,m}$.

The determinant profile itself has a classical finite spectral meaning:
Proposition~\ref{prop:exterior-profile} identifies it, up to fixed
constants, with the exterior singular-value profile of an ordinary
Vandermonde matrix. That identity alone proves no statistical attainment.
Similarly, the arbitrary-prefix Hermite statement in
Lemma~\ref{lem:prefix-hermite} is classical \cite[Proposition 7]{deBoor}.
The attained normalized product, its acquired-history probability and
its compatible transitions are separate mathematical assertions.

The computational step does not diagonalize that geometry. It constructs
a finite net from all grid histories, applies farthest-first selection to
dyadic approximations, and compiles every transition before the stream
begins. The greedy covering mechanism is classical \cite{Gonzalez}.
Theorem~\ref{thm:compact-compiler} states the construction for arbitrary
compact Lipschitz state systems with finite evaluation data, not only for
one hidden scalar. For the present experiments its numerical inputs are
a finite table of prior moments. Positive evidence controls approximation
errors; raw moments avoid singular coordinates. The lower bound is
inherited from the same experiment, not from a changed finite-alphabet
acquisition problem.

This yields a further consequence of the collision law. A high-order
contact between exponents can make the smallest nonzero gap much smaller
than the numerical mesh. Nevertheless the optimal memory order is
attained using a precision set only by $M$, without first identifying
the collision tree. Corollary~\ref{cor:digital-contact} states this for
the intersecting two-parameter arrangement. This is a finite-program
consequence of the attained geometry, not an additional spectral invariant.

The principal argument proceeds from the experiment and exact tangent,
through the isolated analytic inputs, to the intrinsic law and its finite
compilation. Complete proofs for the observation algebra, the affine
specializations, five- and seven-trial laws, and the separate decision
and mechanical comparisons are given in the appendices. They remain
statements with their own hypotheses, rather than claims about arbitrary
rewards or arbitrary hidden-variable models.
